\documentclass[11pt,a4paper]{article}
\usepackage[utf8]{inputenc}
\usepackage[margin=0.8in]{geometry}
\usepackage{listings}
\usepackage{xcolor}
\usepackage{titlesec}
\usepackage{hyperref}
\usepackage{fancyhdr}

\definecolor{codebg}{RGB}{245,245,245}
\definecolor{codegreen}{RGB}{0,128,0}
\definecolor{codegray}{RGB}{128,128,128}
\definecolor{codeblue}{RGB}{0,0,180}
\definecolor{codepurple}{RGB}{128,0,128}

\lstset{
    language=C++,
    backgroundcolor=\color{codebg},
    basicstyle=\ttfamily\footnotesize,
    keywordstyle=\color{codeblue}\bfseries,
    commentstyle=\color{codegreen}\itshape,
    stringstyle=\color{codepurple},
    numberstyle=\tiny\color{codegray},
    numbers=left,
    numbersep=5pt,
    frame=single,
    breaklines=true,
    breakatwhitespace=false,
    tabsize=4,
    showstringspaces=false,
    captionpos=b,
    morekeywords={override,nullptr,string,shared_ptr,unique_ptr,make_shared,make_unique,
                  lock_guard,unique_lock,mutex,thread,vector,unordered_map,map,
                  chrono,steady_clock,time_point,atomic},
}

\pagestyle{fancy}
\fancyhf{}
\fancyhead[L]{\textbf{LLD --- Meeting Room}}
\fancyhead[R]{\thepage}

\title{\Huge\textbf{Meeting Room Booking}\\[0.3em]\Large Low-Level Design --- C++ Concurrency}
\author{LLD Interview Preparation}
\date{}

\begin{document}
\maketitle
\tableofcontents
\newpage

% ============================================================
\section{File Structure}
% ============================================================
\begin{verbatim}
MeetingRoom/
  models.hpp       -- Interval, Meeting, Room, BookingPolicy, Observers
  managers.hpp     -- MeetingMgr (Singleton, header)
  managers.cpp     -- MeetingMgr (implementation)
  main.cpp         -- Demo: policies, concurrent booking, cancellation
\end{verbatim}

\textbf{Compilation:}
\begin{verbatim}
g++ -std=c++17 -pthread -o meeting_room.exe main.cpp managers.cpp
\end{verbatim}

\subsection{Concurrency Summary}
\begin{itemize}
    \item \textbf{Level 1:} \texttt{operationMtx} (MeetingMgr) --- rooms map, policy, observers
    \item \textbf{Level 2:} \texttt{roomMtx} (per-Room) --- room calendar
    \item \textbf{lock\_guard:} \texttt{addRoom}, \texttt{setPolicy}, \texttt{addObserver}, \texttt{listAllRooms}
    \item \textbf{unique\_lock:} \texttt{bookRoom} (release before room-level lock)
    \item \textbf{Snapshot pattern:} observers notified outside lock
\end{itemize}

\subsection{Design Patterns}
\begin{itemize}
    \item \textbf{Singleton:} MeetingMgr
    \item \textbf{Strategy:} BookingPolicy (Standard, MaxDuration, AdvanceBooking)
    \item \textbf{Observer:} EmailNotifier, SlackNotifier, SMSNotifier
\end{itemize}

\subsection{Smart Pointers}
\begin{itemize}
    \item \texttt{shared\_ptr<Room>} --- map + booking threads hold ref
    \item \texttt{unique\_ptr<BookingPolicy>} --- manager sole owner
    \item \texttt{shared\_ptr<NotificationObserver>} --- manager + caller
\end{itemize}

% ============================================================
\newpage
\section{models.hpp}
% ============================================================
\begin{lstlisting}[title={\texttt{models.hpp}}]
#ifndef MODELS_HPP
#define MODELS_HPP
#include<iostream> 
#include<vector>
#include<mutex>
#include<algorithm>
#include<string>
using namespace std;


class Interval {
    public:
    int startTime;
    int endTime;

    Interval(int start, int end) : startTime(start), endTime(end) {
        if(start >= end){
            throw invalid_argument("Start time must be before end Time");
        }
    }

    bool overlaps(const Interval& other) const {
        return this->startTime < other.endTime && other.startTime < this->endTime;
    }
};

class Meeting {
    public:
    string id;
    Interval interval;
    string userId;

    Meeting(const string& id, const Interval& interval, const string& userId) 
       : id(id), interval(interval), userId(userId) {}

};


//handles its own concurrency
class Room{
    string id;
    int capacity;
    vector<Meeting> calender;
    mutable mutex roomMtx;

    public:

    Room(const string& id, const int& capacity)
       : id(id), capacity(capacity) {}

    string getId() const {return id;}   
    int getCapacity() const {return capacity;}

    bool book(const string& meetingId, const Interval& interval, const string& userId){
        // 1. Lock this specific room
        lock_guard<mutex> lock(roomMtx);

        // 2. Critical Section: Check for overlaps
        for(const auto& meeting : calender){
            if(meeting.interval.overlaps(interval)){
                return false; //conflict found
            }
        }

        // 3. No conflict: Add meeting
        calender.emplace_back(meetingId, interval, userId);
        return true;
        // 4. Lock automatically releases here
    }
    
    bool cancelBooking(const string& meetingId) {
        lock_guard<mutex> lock(roomMtx);
        
        auto it = find_if(calender.begin(), calender.end(),
            [&meetingId](const Meeting& m) { return m.id == meetingId; });
        
        if (it != calender.end()) {
            calender.erase(it);
            return true;
        }
        return false;
    }
    
    bool isAvailable(int start, int end) const {
        lock_guard<mutex> lock(roomMtx);
        Interval requestedInterval(start, end);
        
        for (const auto& meeting : calender) {
            if (meeting.interval.overlaps(requestedInterval)) {
                return false;
            }
        }
        return true;
    }
    
    void displayBookings() const {
        lock_guard<mutex> lock(roomMtx);
        cout << "Room " << id << " (Capacity: " << capacity << ") Bookings:\n";
        
        if (calender.empty()) {
            cout << "  No bookings.\n";
        } else {
            for (const auto& meeting : calender) {
                cout << "  [" << meeting.interval.startTime << "-" 
                     << meeting.interval.endTime << "] " 
                     << meeting.userId << " (ID: " << meeting.id << ")\n";
            }
        }
    }
    
    int getBookingCount() const {
        lock_guard<mutex> lock(roomMtx);
        return calender.size();
    }
};


// ========== STRATEGY PATTERN ==========
// Defines different booking policies

class BookingPolicy {
public:
    virtual ~BookingPolicy() = default;
    virtual bool canBook(const Room* room, const Interval& interval, const string& userId) const = 0;
    virtual string getPolicyName() const = 0;
};

class StandardBookingPolicy : public BookingPolicy {
public:
    bool canBook(const Room* room, const Interval& interval, const string& userId) const override {
        return true;  // Anyone can book
    }
    
    string getPolicyName() const override {
        return "Standard (No restrictions)";
    }
};

class MaxDurationPolicy : public BookingPolicy {
    int maxDuration;
public:
    MaxDurationPolicy(int maxHours) : maxDuration(maxHours) {}
    
    bool canBook(const Room* room, const Interval& interval, const string& userId) const override {
        int duration = interval.endTime - interval.startTime;
        if (duration > maxDuration) {
            cout << "Policy Violation: Max booking duration is " << maxDuration << " hours.\n";
            return false;
        }
        return true;
    }
    
    string getPolicyName() const override {
        return "Max Duration (" + to_string(maxDuration) + " hours)";
    }
};

class AdvanceBookingPolicy : public BookingPolicy {
    int minAdvanceHours;
public:
    AdvanceBookingPolicy(int hours) : minAdvanceHours(hours) {}
    
    bool canBook(const Room* room, const Interval& interval, const string& userId) const override {
        // In real system, would compare with current time
        // Simplified: check if booking is not immediate
        if (interval.startTime < minAdvanceHours) {
            cout << "Policy Violation: Must book at least " << minAdvanceHours << " hours in advance.\n";
            return false;
        }
        return true;
    }
    
    string getPolicyName() const override {
        return "Advance Booking (" + to_string(minAdvanceHours) + " hours notice)";
    }
};


// ========== OBSERVER PATTERN ==========
// Notification system for booking events

class NotificationObserver {
public:
    virtual ~NotificationObserver() = default;
    virtual void update(const string& message) = 0;
    virtual string getObserverName() const = 0;
};

class EmailNotifier : public NotificationObserver {
    string email;
public:
    EmailNotifier(const string& email) : email(email) {}
    
    void update(const string& message) override {
        cout << "[EMAIL to " << email << "]: " << message << endl;
    }
    
    string getObserverName() const override {
        return "EmailNotifier(" + email + ")";
    }
};

class SlackNotifier : public NotificationObserver {
    string channel;
public:
    SlackNotifier(const string& channel) : channel(channel) {}
    
    void update(const string& message) override {
        cout << "[SLACK #" << channel << "]: " << message << endl;
    }
    
    string getObserverName() const override {
        return "SlackNotifier(#" + channel + ")";
    }
};

class SMSNotifier : public NotificationObserver {
    string phoneNumber;
public:
    SMSNotifier(const string& phone) : phoneNumber(phone) {}
    
    void update(const string& message) override {
        cout << "[SMS to " << phoneNumber << "]: " << message << endl;
    }
    
    string getObserverName() const override {
        return "SMSNotifier(" + phoneNumber + ")";
    }
};


#endif // MODELS_HPP
\end{lstlisting}

% ============================================================
\newpage
\section{managers.hpp}
% ============================================================
\begin{lstlisting}[title={\texttt{managers.hpp}}]
#ifndef MANAGERS_HPP
#define MANAGERS_HPP

#include<iostream>
#include<mutex>
#include<unordered_map>
#include<memory>
#include<thread>
#include "models.hpp"
using namespace std;

// ============================================================
// LOCKING STRATEGY: PESSIMISTIC vs OPTIMISTIC
// ----------------------------------------------------------
// PESSIMISTIC LOCKING (what we use here):
//   Lock BEFORE reading/writing shared state. Assumes conflicts
//   are likely. Simple, safe, but can reduce throughput.
//   e.g., lock_guard on the rooms map before looking up a room.
//
// OPTIMISTIC LOCKING:
//   Read without lock, do work, then check-and-commit atomically.
//   Uses version numbers or CAS (compare-and-swap).
//   Better throughput when conflicts are rare.
//   e.g., Room could have a version field: read version, prepare
//   booking, then CAS(old_version, new_version) -- retry if changed.
//
// WHY PESSIMISTIC HERE?
//   Meeting rooms have high contention (many users booking same rooms).
//   Optimistic would cause lots of retries. Pessimistic is simpler
//   and correct for this use case.
// ============================================================

// ============================================================
// SMART POINTER CHOICES:
// ----------------------------------------------------------
// shared_ptr<Room>     -> Room is in the map AND may be held by
//                         threads doing bookings. If room is removed
//                         mid-booking, shared_ptr keeps it alive.
// unique_ptr<Policy>   -> Manager is the SOLE owner of the policy.
//                         Nobody else holds a reference. Transferred
//                         via move(). Lightweight, no ref counting.
// shared_ptr<Observer> -> Observers are shared: caller keeps a
//                         reference (to remove later), manager's
//                         vector also holds one. Shared ownership.
// raw ptr              -> Not used here. Could be used for non-owning
//                         observer references if we didn't need
//                         shared lifetime management.
// ============================================================

// ============================================================
// LOCK CHOICES:
// ----------------------------------------------------------
// lock_guard    -> Used for simple scoped locks: addRoom, removeRoom
//                  lookup, setPolicy, add/removeObserver.
//                  Holds lock for entire scope. No early release.
//
// unique_lock   -> Used in bookRoom: we lock to find the room and
//                  read the policy, then UNLOCK before the actual
//                  room->book() call (which has its own per-room lock).
//                  This avoids holding the manager lock during the
//                  booking operation (fine-grained locking).
//                  Also used when we need to release lock before
//                  notifying observers (avoids deadlock).
// ============================================================

class MeetingMgr{
    private:
    static mutex singletonMtx;              // only for singleton creation

    // operationMtx: guards rooms map, bookingPolicy, observers, etc.
    // Separate from singletonMtx to avoid contention on getInstance()
    mutable mutex operationMtx;

    // shared_ptr<Room>: room lives in map AND may be held by booking
    // threads concurrently. Shared ownership keeps it alive.
    unordered_map<string, shared_ptr<Room>> rooms;
    
    // unique_ptr<BookingPolicy>: sole ownership by manager.
    // Transferred via move(). Only one policy active at a time.
    unique_ptr<BookingPolicy> bookingPolicy;
    
    // shared_ptr<Observer>: shared between manager's vector and
    // the caller (who needs a reference to removeObserver later).
    vector<shared_ptr<NotificationObserver>> observers;

    MeetingMgr();

    // Notify MUST be called WITHOUT operationMtx held (avoids deadlock
    // if an observer callback triggers another manager operation)
    void notifyObservers(const string& message);

    public:
    static shared_ptr<MeetingMgr> getInstance();
    
    // Room Management
    void addRoom(const string& roomId, int capacity);
    bool removeRoom(const string& roomId);
    
    // Booking Operations
    bool bookRoom(const string& roomId, int start, int end, const string& userId);
    bool cancelBooking(const string& roomId, const string& meetingId);
    
    // Query Operations
    void viewRoomBookings(const string& roomId) const;
    void listAllRooms() const;
    vector<string> findAvailableRooms(int start, int end) const;
    
    // Strategy Pattern Methods
    void setBookingPolicy(unique_ptr<BookingPolicy> policy);
    
    // Observer Pattern Methods
    void addObserver(shared_ptr<NotificationObserver> observer);
    void removeObserver(shared_ptr<NotificationObserver> observer);

    MeetingMgr(const MeetingMgr&) = delete;
    void operator=(const MeetingMgr&) = delete;
};


#endif
\end{lstlisting}

% ============================================================
\newpage
\section{managers.cpp}
% ============================================================
\begin{lstlisting}[title={\texttt{managers.cpp}}]
#include<iostream>
#include <string>
#include "managers.hpp"
#include "models.hpp"

using namespace std;

mutex MeetingMgr::singletonMtx;

MeetingMgr::MeetingMgr() {
    bookingPolicy = make_unique<StandardBookingPolicy>();
}

// Thread-safe singleton using function-local static
// The C++11 standard guarantees static local init is thread-safe
shared_ptr<MeetingMgr> MeetingMgr::getInstance(){
    lock_guard<mutex> lock(singletonMtx);
    static shared_ptr<MeetingMgr> instance(new MeetingMgr());
    return instance;
}

// ============================================================
// PESSIMISTIC LOCKING PATTERN used throughout:
// 1. Acquire lock (operationMtx)
// 2. Read/write shared state
// 3. Release lock BEFORE I/O (notifications, printing)
// This prevents deadlocks and reduces lock hold time.
// ============================================================

void MeetingMgr::addRoom(const string& roomId, int capacity){
    // lock_guard: simple scope, hold lock while modifying rooms map
    {
        lock_guard<mutex> lock(operationMtx);
        rooms[roomId] = make_shared<Room>(roomId, capacity);
    }
    // Lock released BEFORE I/O and notification (avoids deadlock)
    cout << "[Thread " << this_thread::get_id() << "] "
         << "Room " << roomId << " added (Capacity: " << capacity << ")" << endl;
    notifyObservers("Room " + roomId + " has been added to the system.");
}

bool MeetingMgr::removeRoom(const string& roomId) {
    shared_ptr<Room> room;
    
    {
        // lock_guard: simple lookup + erase, hold for entire block
        lock_guard<mutex> lock(operationMtx);
        auto it = rooms.find(roomId);
        
        if (it == rooms.end()) {
            cout << "Error: Room " << roomId << " does not exist." << endl;
            return false;
        }
        
        // shared_ptr ref count keeps room alive even after map erase.
        // If another thread is mid-booking on this room, the room
        // object stays alive until that thread's shared_ptr drops.
        room = it->second;  // ref count: map(1) + local(1) = 2
        rooms.erase(it);    // ref count: local(1) = 1
    }
    // Lock released BEFORE notification
    
    if (room.use_count() > 1) {
        cout << "[Thread " << this_thread::get_id() << "] "
             << "Room " << roomId << " removed but still in use by " 
             << (room.use_count() - 1) << " active operation(s). "
             << "Will be deleted when all operations complete." << endl;
    } else {
        cout << "[Thread " << this_thread::get_id() << "] "
             << "Room " << roomId << " successfully removed." << endl;
    }
    
    notifyObservers("Room " + roomId + " has been removed from the system.");
    return true;
}

bool MeetingMgr::bookRoom(const string& roomId, int start, int end, const string& userId){
    shared_ptr<Room> room;
    
    // ============================================================
    // FINE-GRAINED LOCKING (2-level pessimistic):
    // Level 1: operationMtx (manager-level) -> protects rooms map + policy
    // Level 2: roomMtx (per-room) -> protects that room's calendar
    //
    // unique_lock here because we RELEASE operationMtx after
    // looking up the room and checking the policy. Then room->book()
    // only takes the per-room lock. This way, bookings to DIFFERENT
    // rooms can proceed in parallel.
    // ============================================================
    unique_lock<mutex> lock(operationMtx);
    
    auto it = rooms.find(roomId);
    if (it == rooms.end()) {
        cout << "Error: Room " << roomId << " does not exist." << endl;
        return false;
    }
    room = it->second;  // shared_ptr: room stays alive even if removed from map

    try {
        Interval interval(start, end);
        
        // Check policy while we still hold operationMtx
        // (policy could be swapped by setBookingPolicy on another thread)
        if (!bookingPolicy->canBook(room.get(), interval, userId)) {
            cout << "POLICY REJECTION: " << userId << " cannot book Room " << roomId 
                 << " [" << start << "-" << end << "]" << endl;
            return false;
        }
        
        // Release manager lock BEFORE room-level operation
        // This is the KEY benefit of unique_lock over lock_guard:
        // Other threads can now book DIFFERENT rooms in parallel
        lock.unlock();
        
        string meetingId = userId + "_" + to_string(start);
        
        // room->book() acquires its own per-room roomMtx (Level 2 lock)
        // This is pessimistic locking at the room level
        bool success = room->book(meetingId, interval, userId);
        
        if (success) {
            cout << "[Thread " << this_thread::get_id() << "] "
                 << "SUCCESS: " << userId << " booked Room " << roomId 
                 << " [" << start << "-" << end << "]" << endl;
            // Notify OUTSIDE all locks (avoids deadlock)
            notifyObservers(userId + " booked Room " + roomId + " [" + 
                          to_string(start) + "-" + to_string(end) + "]");
            return true;
        } else {
            cout << "[Thread " << this_thread::get_id() << "] "
                 << "FAIL: " << userId << " could not book Room " << roomId 
                 << " [" << start << "-" << end << "] (Conflict)" << endl;
            return false;
        }
    } catch (const exception& e) {
        cout << "Error: " << e.what() << endl;
        return false;
    }
}

bool MeetingMgr::cancelBooking(const string& roomId, const string& meetingId) {
    shared_ptr<Room> room;
    
    {
        // lock_guard: simple lookup, hold for short duration
        lock_guard<mutex> lock(operationMtx);
        auto it = rooms.find(roomId);
        if (it == rooms.end()) {
            cout << "Error: Room " << roomId << " does not exist." << endl;
            return false;
        }
        room = it->second;
    }
    // Manager lock released. room->cancelBooking uses its own roomMtx.
    
    bool success = room->cancelBooking(meetingId);
    
    if (success) {
        cout << "[Thread " << this_thread::get_id() << "] "
             << "Booking " << meetingId << " cancelled in Room " << roomId << endl;
        notifyObservers("Booking " + meetingId + " cancelled in Room " + roomId);
        return true;
    } else {
        cout << "Booking " << meetingId << " not found in Room " << roomId << endl;
        return false;
    }
}

void MeetingMgr::viewRoomBookings(const string& roomId) const {
    shared_ptr<Room> room;
    
    {
        lock_guard<mutex> lock(operationMtx);
        auto it = rooms.find(roomId);
        if (it == rooms.end()) {
            cout << "Error: Room " << roomId << " does not exist." << endl;
            return;
        }
        room = it->second;
    }
    // Manager lock released. room->displayBookings uses its own roomMtx.
    room->displayBookings();
}

void MeetingMgr::listAllRooms() const {
    lock_guard<mutex> lock(operationMtx);
    
    cout << "\n=== All Meeting Rooms ===" << endl;
    if (rooms.empty()) {
        cout << "No rooms available." << endl;
        return;
    }
    
    for (const auto& [roomId, room] : rooms) {
        cout << "  " << roomId << " (Capacity: " << room->getCapacity() 
             << ", Bookings: " << room->getBookingCount() << ")" << endl;
    }
}

vector<string> MeetingMgr::findAvailableRooms(int start, int end) const {
    lock_guard<mutex> lock(operationMtx);
    vector<string> availableRooms;
    
    for (const auto& [roomId, room] : rooms) {
        if (room->isAvailable(start, end)) {
            availableRooms.push_back(roomId);
        }
    }
    
    return availableRooms;
}

void MeetingMgr::setBookingPolicy(unique_ptr<BookingPolicy> policy) {
    // lock_guard: simple one-liner write to shared state
    lock_guard<mutex> lock(operationMtx);
    bookingPolicy = move(policy);
    cout << "Policy changed to: " << bookingPolicy->getPolicyName() << endl;
}

void MeetingMgr::addObserver(shared_ptr<NotificationObserver> observer) {
    lock_guard<mutex> lock(operationMtx);
    observers.push_back(observer);
    cout << "Observer added: " << observer->getObserverName() << endl;
}

void MeetingMgr::removeObserver(shared_ptr<NotificationObserver> observer) {
    lock_guard<mutex> lock(operationMtx);
    auto it = find(observers.begin(), observers.end(), observer);
    if (it != observers.end()) {
        cout << "Observer removed: " << (*it)->getObserverName() << endl;
        observers.erase(it);
    }
}

// CRITICAL: This is called WITHOUT operationMtx held.
// We take a snapshot of observers under lock, then notify outside lock.
// This prevents deadlock if an observer callback triggers another
// manager operation (which would try to acquire operationMtx again).
void MeetingMgr::notifyObservers(const string& message) {
    vector<shared_ptr<NotificationObserver>> snapshot;
    {
        lock_guard<mutex> lock(operationMtx);
        snapshot = observers;  // copy the vector (shared_ptrs are cheap to copy)
    }
    // Notify outside the lock -- safe because we hold shared_ptrs
    for (const auto& observer : snapshot) {
        observer->update(message);
    }
}
\end{lstlisting}

% ============================================================
\newpage
\section{main.cpp}
% ============================================================
\begin{lstlisting}[title={\texttt{main.cpp}}]
#include<iostream>
#include<thread>
#include<chrono>
#include "managers.hpp"
using namespace std;


void userSimulation(string userId, string roomId, int start, int end) {
    auto mgr = MeetingMgr::getInstance();
    mgr->bookRoom(roomId, start, end, userId);
}

void demonstrateRemoveRoom() {
    auto mgr = MeetingMgr::getInstance();
    
    cout << "\n=== Demonstrating Room Removal with Active Usage ===" << endl;
    
    mgr->addRoom("R3", 6);
    
    // Thread tries to book R3 while main thread removes it
    // shared_ptr<Room> keeps the room alive even after map erase
    thread t1([&]() {
        cout << "Thread 1: Got reference to R3" << endl;
        this_thread::sleep_for(chrono::milliseconds(100));
        mgr->bookRoom("R3", 14, 15, "UserE");
    });
    
    this_thread::sleep_for(chrono::milliseconds(50));
    mgr->removeRoom("R3");
    
    t1.join();
}

int main(){
    auto mgr = MeetingMgr::getInstance();

    cout << "=== Meeting Room Booking System ===" << endl << endl;
    
    // ========== 1. OBSERVER PATTERN DEMO ==========
    cout << "--- Setting up Notification Observers ---" << endl;
    // shared_ptr<Observer>: shared between caller (to remove later)
    // and manager's observer vector. Both need ownership.
    auto emailNotifier = make_shared<EmailNotifier>("admin@company.com");
    auto slackNotifier = make_shared<SlackNotifier>("meeting-rooms");
    auto smsNotifier = make_shared<SMSNotifier>("+1-555-0123");
    
    mgr->addObserver(emailNotifier);
    mgr->addObserver(slackNotifier);
    mgr->addObserver(smsNotifier);
    
    
    // ========== 2. ROOM SETUP ==========
    cout << "\n--- Adding Rooms ---" << endl;
    mgr->addRoom("R1", 10);
    mgr->addRoom("R2", 4);
    
    mgr->listAllRooms();
    
    
    // ========== 3. STRATEGY PATTERN DEMO ==========
    cout << "\n--- Testing Booking Policies (Strategy Pattern) ---" << endl;
    
    // unique_ptr<BookingPolicy>: manager is sole owner. Transferred via move().
    // When we setBookingPolicy, old policy is auto-deleted, new one takes over.
    cout << "\nUsing Standard Policy:" << endl;
    mgr->bookRoom("R1", 9, 14, "Alice");
    
    cout << "\nChanging to Max Duration Policy (3 hours max):" << endl;
    mgr->setBookingPolicy(make_unique<MaxDurationPolicy>(3));
    mgr->bookRoom("R1", 15, 20, "Bob");     // 5 hours - should FAIL
    mgr->bookRoom("R1", 15, 17, "Bob");     // 2 hours - should work
    
    cout << "\nChanging to Advance Booking Policy (min 5 hours notice):" << endl;
    mgr->setBookingPolicy(make_unique<AdvanceBookingPolicy>(5));
    mgr->bookRoom("R2", 3, 4, "Charlie");     // 3hrs notice - should FAIL
    mgr->bookRoom("R2", 10, 11, "Charlie");   // 10hrs notice - should work
    
    
    // ========== 4. CONCURRENT BOOKING DEMO ==========
    // ============================================================
    // PESSIMISTIC LOCKING IN ACTION:
    // - bookRoom uses unique_lock on operationMtx to find room + check policy
    // - Then RELEASES operationMtx (fine-grained: other rooms can proceed)
    // - room->book() uses its own lock_guard on roomMtx (per-room lock)
    // - Two threads booking SAME room at SAME time: one wins, one gets conflict
    // - Two threads booking DIFFERENT rooms: proceed in parallel (no contention)
    // ============================================================
    cout << "\n--- Concurrent Booking (Pessimistic Locking) ---" << endl;
    cout << "UserA & UserB race for same slot on R1 (one will fail)." << endl;
    cout << "UserC books different time on R1 (no conflict)." << endl;
    cout << "UserD books R2 in parallel (different room, no contention)." << endl << endl;
    
    mgr->setBookingPolicy(make_unique<StandardBookingPolicy>());
    
    thread t1(userSimulation, "UserA", "R1", 20, 21);
    thread t2(userSimulation, "UserB", "R1", 20, 21); // Conflict!
    thread t3(userSimulation, "UserC", "R1", 21, 22); // Different time
    thread t4(userSimulation, "UserD", "R2", 20, 21); // Different room
    
    t1.join();
    t2.join();
    t3.join();
    t4.join();
    
    
    // ========== 5. QUERY OPERATIONS ==========
    cout << "\n--- Viewing Room Bookings ---" << endl;
    mgr->viewRoomBookings("R1");
    mgr->viewRoomBookings("R2");
    
    cout << "\n--- Finding Available Rooms ---" << endl;
    auto available = mgr->findAvailableRooms(20, 21);
    cout << "Rooms available [20-21]: ";
    if (available.empty()) {
        cout << "None" << endl;
    } else {
        for (const auto& room : available) {
            cout << room << " ";
        }
        cout << endl;
    }
    
    
    // ========== 6. CANCEL BOOKING ==========
    cout << "\n--- Cancelling Booking ---" << endl;
    mgr->cancelBooking("R1", "Alice_9");
    mgr->viewRoomBookings("R1");
    
    
    // ========== 7. ROOM REMOVAL DEMO ==========
    // shared_ptr<Room> keeps room alive if another thread holds a ref
    cout << "\n--- Removing Room ---" << endl;
    mgr->removeRoom("R2");
    mgr->listAllRooms();
    
    mgr->bookRoom("R2", 25, 26, "UserX");  // Should fail - room removed
    
    
    // ========== 8. ADVANCED: Remove room while in use ==========
    demonstrateRemoveRoom();
    
    
    // ========== 9. OBSERVER MANAGEMENT ==========
    cout << "\n--- Removing Observer ---" << endl;
    mgr->removeObserver(smsNotifier);
    
    cout << "\n--- Final Booking (fewer notifications) ---" << endl;
    mgr->bookRoom("R1", 30, 31, "FinalUser");
    
    // ============================================================
    // SUMMARY: CONCURRENCY ARCHITECTURE
    // ----------------------------------------------------------
    // Level 1: operationMtx (manager) -> rooms map, policy, observers
    //   - lock_guard for simple reads/writes (addRoom, setPolicy, etc.)
    //   - unique_lock in bookRoom (release before room-level lock)
    //
    // Level 2: roomMtx (per-room) -> that room's calendar
    //   - lock_guard always (simple check+book is atomic)
    //
    // DEADLOCK PREVENTION:
    //   - Always acquire Level 1 before Level 2 (never reverse)
    //   - Release Level 1 before calling notifyObservers
    //   - notifyObservers takes a snapshot under lock, notifies outside
    //
    // WHY PESSIMISTIC?
    //   Meeting rooms = high contention. Pessimistic is simpler
    //   and avoids retry storms that optimistic would cause.
    // ============================================================
    
    cout << "\n=== Simulation Complete ===" << endl;
    return 0;
}
\end{lstlisting}

\end{document}
